\section*{Capítulo \ref{ch:g10:lines} --- Ecuaciones de rectas y sistemas lineales}

\begin{solution}{exo:g10:lines:1}
$y = 2x - 3$: \omterm{def:g10:reffunc:affine}{pendiente} $2$, \omterm{def:g10:reffunc:affine}{ordenada en el origen} $-3$.
$y = -\frac13 x + 1$: \omterm{def:g10:reffunc:affine}{pendiente} $-\frac13$, \omterm{def:g10:reffunc:affine}{ordenada en el origen} $1$.
$y = 4$: \omterm{def:g10:reffunc:affine}{pendiente} $0$, \omterm{def:g10:reffunc:affine}{ordenada en el origen} $4$ (recta horizontal).
$x = -2$: recta vertical, sin \omterm{def:g10:reffunc:affine}{pendiente} y sin \omterm{thm:g10:lines:reduced}{ecuación explícita}.
\end{solution}

\begin{solution}{exo:g10:lines:2}
$3 \times 2 + 1 = 7$: sí, $A$ está sobre la recta.
$3 \times (-1) + 1 = -2 \neq -1$: $B$ no lo está.
$3 \times 4 + 1 = 13$: sí, $C$ está sobre la recta.
\end{solution}

\begin{solution}{exo:g10:lines:3}
(a) $m = \frac{8 - 2}{3 - 0} = 2$ y $p = y_A = 2$ (el punto $A$ está sobre
el eje $y$): $y = 2x + 2$.

(b) $m = \frac{-4 - 5}{2 - (-1)} = \frac{-9}{3} = -3$; después
$5 = -3 \times (-1) + p$ da $p = 2$: $y = -3x + 2$.

(c) $x_A = x_B = 4$: recta vertical $x = 4$.
\end{solution}

\begin{solution}{exo:g10:lines:4}
$y = \frac{6x+2}{2} = 3x + 1$. Las rectas de \omterm{def:g10:reffunc:affine}{pendiente} $3$ son
$y = 3x - 1$, $y = 3x + 5$ e $y = 3x + 1$: esas tres son paralelas entre
sí; ninguna pareja de ellas coincide (las \omterm{def:g10:reffunc:affine}{ordenadas en el origen} son
distintas) y ninguna es paralela a $y = -3x + 2$.
\end{solution}

\begin{solution}{exo:g10:lines:5}
Primer sistema: sustituimos $y = 2x - 1$ en $3x + y = 9$:
$3x + 2x - 1 = 9$, luego $5x = 10$, $x = 2$ y después $y = 3$. Solución:
$(2, 3)$.

Segundo sistema: de $x - y = 4$, $x = y + 4$. Entonces
$2(y + 4) + 3y = 3$, luego $5y = -5$, $y = -1$ y después $x = 3$.
Solución: $(3, -1)$.
\end{solution}

\begin{solution}{exo:g10:lines:6}
Primer sistema: sumando las ecuaciones se elimina $y$: $2x = 14$, $x = 7$;
después $7 + y = 10$ da $y = 3$. Solución: $(7, 3)$.

Segundo sistema: multiplicamos la primera \omterm{def:g10:algebra:equation}{ecuación} por $5$ y la segunda
por $-2$:
\[
\begin{cases}
15x + 10y = 60 \\
-4x - 10y = -16
\end{cases}
\]
Sumando: $11x = 44$, luego $x = 4$; después $3 \times 4 + 2y = 12$ da
$y = 0$. Solución: $(4, 0)$.
\end{solution}

\begin{solution}{exo:g10:lines:7}
Calculamos $ab' - a'b$ en cada caso.

Primero: $2 \times 2 - (-4) \times (-1) = 4 - 4 = 0$, y la segunda
\omterm{def:g10:algebra:equation}{ecuación} es $-2$ veces la primera: dos veces la misma recta, infinitas
soluciones.

Segundo: de nuevo $ab' - a'b = 0$, pero los segundos miembros no están en
la razón $-2$ ($1 \neq -6$): dos rectas paralelas distintas, sin solución.

Tercero: $2 \times 1 - 1 \times (-1) = 3 \neq 0$: exactamente una
solución.
\end{solution}

\begin{solution}{exo:g10:lines:8}
Sea $x$ el número de entradas infantiles e $y$ el de entradas de adulto:
\[
\begin{cases}
x + y = 200 \\
5x + 9y = 1352 .
\end{cases}
\]
De la primera \omterm{def:g10:algebra:equation}{ecuación}, $x = 200 - y$; sustituyendo,
$5(200 - y) + 9y = 1352$, luego $1000 + 4y = 1352$, $y = 88$ y después
$x = 112$. Así pues, $112$ entradas infantiles y $88$ de adulto.
Comprobación: $5 \times 112 + 9 \times 88 = 560 + 792 = 1352$.
\end{solution}

\begin{solution}{exo:g10:lines:9}
\emph{1.} Paralela significa misma \omterm{def:g10:reffunc:affine}{pendiente} $2$: $y = 2x + p$ con
$4 = 2 \times 1 + p$, luego $p = 2$: $d'\colon y = 2x + 2$.

\emph{2.} Sobre el eje $x$, $y = 0$: $2x + 2 = 0$ da $x = -1$. El punto de
corte es $(-1, 0)$.
\end{solution}

\begin{solution}{exo:g10:lines:10}
\emph{1.} \omterm{prop:g10:coordgeom:midpoint}{Punto medio} de $[BC]$: $A' = \left(\frac{6+2}{2},
\frac{2+4}{2}\right) = (4, 3)$. La mediana desde $A(0,0)$ tiene \omterm{def:g10:reffunc:affine}{pendiente}
$\frac{3 - 0}{4 - 0} = \frac34$: \omterm{def:g10:algebra:equation}{ecuación} $y = \frac34 x$.

\omterm{prop:g10:coordgeom:midpoint}{Punto medio} de $[AC]$: $B' = (1, 2)$. La mediana desde $B(6, 2)$ tiene
\omterm{def:g10:reffunc:affine}{pendiente} $\frac{2 - 2}{1 - 6} = 0$: es la recta horizontal $y = 2$.

\emph{2.} Corte: $\frac34 x = 2$ da $x = \frac83$, luego
$G\left(\frac83, 2\right)$. La tercera mediana une $C(2,4)$ con el \omterm{prop:g10:coordgeom:midpoint}{punto
medio} de $[AB]$, $C' = (3, 1)$; su \omterm{def:g10:reffunc:affine}{pendiente} es
$\frac{1 - 4}{3 - 2} = -3$, \omterm{def:g10:algebra:equation}{ecuación} $y = -3x + 10$. En $x = \frac83$:
$y = -8 + 10 = 2$: $G$ está sobre ella. Y, en efecto,
$G = \left(\frac{0 + 6 + 2}{3}, \frac{0 + 2 + 4}{3}\right)$: el baricentro
promedia las \omterm{def:g10:coordgeom:system}{coordenadas} de los vértices.
\end{solution}

\begin{solution}{pb:g10:lines:1}
\textbf{1.} \omterm{def:g10:reffunc:affine}{Pendiente} $\frac{8 - 2}{3 - 1} = 3$; pasando por $(1, 2)$:
$y = 3x - 1$.

\textbf{2.} $y = 3x - 1$ e $y = 3x + 4$: misma \omterm{def:g10:reffunc:affine}{pendiente}, paralelas (y
distintas). Cortando $y = 3x - 1$ con $y = -x + 7$: $3x - 1 = -x + 7$ da
$x = 2$, $y = 5$: el punto $(2, 5)$.

\textbf{3.} Los dos puntos comparten $x = 2$: la recta es vertical, de
\omterm{def:g10:algebra:equation}{ecuación} $x = 2$. Una forma explícita $y = mx + p$ responde a ``un $y$ por
cada $x$'', cosa que una recta vertical se niega a hacer: su \omterm{def:g10:reffunc:affine}{pendiente}
sería infinita.

\textbf{4.} $3 \times 4 - 1 = 11$: sí, $(4, 11)$ está sobre la recta.
$3 \times (-1) - 1 = -4 \neq -5$: no.

\textbf{5.} $y = -2x + 3$; corta a los ejes en $(0, 3)$ y
$\left(\frac32, 0\right)$. Área del triángulo:
$\frac12 \times \frac32 \times 3 = \frac94 = 2.25$.

\textbf{6.} Sustituyendo: $3x + 2x - 3 = 7$, $x = 2$, $y = 1$: solución
$(2, 1)$.

\textbf{7.} Sumando: $7x = 7$, $x = 1$; después $3y = 6$, $y = 2$:
solución $(1, 2)$.

\textbf{8.} Primer sistema: la segunda \omterm{def:g10:algebra:equation}{ecuación} es el doble de la primera:
una recta contada dos veces, con infinitas soluciones (todos los puntos de
$2x - y = 3$). Segundo sistema: los primeros miembros son proporcionales,
pero los segundos no ($2 \times 3 = 6 \neq 1$): dos rectas paralelas, sin
solución. Tricotomía: \omterm{def:g10:reffunc:affine}{pendientes} distintas $\to$ un corte; \omterm{def:g10:reffunc:affine}{pendientes}
iguales y \omterm{def:g10:reffunc:affine}{ordenadas en el origen} distintas $\to$ paralelas, ninguna
solución; todo igual $\to$ la misma recta, infinitas soluciones.

\textbf{9.} Los \omterm{def:g10:vectors:vector}{vectores} directores de las rectas son $(-b, a)$ y
$(-d, c)$ (o bien las \omterm{def:g10:reffunc:affine}{pendientes} $-\frac ab$ y $-\frac cd$), \omterm{def:g10:vectors:collinear}{colineales}
exactamente cuando $ad - bc = 0$: el determinante del
\cref{pb:g10:vectors:1} en un nuevo oficio. Valores: pregunta 7:
$2 \times (-3) - 3 \times 5 = -21 \neq 0$: solución única. Pregunta 8:
$2 \times (-2) - (-1) \times 4 = 0$ en los dos casos:
paralelas-o-idénticas, y son los términos independientes los que separan
los dos casos.

\textbf{10.} $ad - bc = 6 - 2m$: se anula para $m = 3$. Ahí,
$x + 3y = 3$ frente a $2x + 6y = 7$: doblando la primera queda
$2x + 6y = 6 \neq 7$, luego son paralelas y \emph{no hay solución}. Para
$m \neq 3$, la reducción da la solución única $y = \frac{1}{6 - 2m}$ y
después $x = 3 - \frac{m}{6 - 2m}$.

\textbf{11.} $f + c = 35$, $2f + 4c = 94$; dividiendo entre dos,
$f + 2c = 47$, luego $c = 12$ y $f = 23$: veintitrés faisanes y doce
conejos, la misma respuesta que dan los \emph{Nueve capítulos}.

\textbf{12.} $x + y = 30$ y $0.6x + 0.2y = 0.35 \times 30 = 10.5$;
restando $0.2 \times$ la primera, $0.4x = 4.5$: $x = 11.25$ L de jarabe e
$y = 18.75$ L de bebida.

\textbf{13.} Cifras $a$ y $b$: $a + b = 12$ y
$(10a + b) - (10b + a) = 9(a - b) = 18$, luego $a - b = 2$: $a = 7$,
$b = 5$; el número es $75$.

\textbf{14.} Sumando: $5c + 5r = 18$, luego $c + r = 3.60$: un café más un
cruasán. Restando: $c - r = 0.20$. De ahí, $c = 1.90$ y $r = 1.70$ euros:
el atajo simétrico lo ha resuelto sin ninguna sustitución.

\textbf{15.} Doblando el primer anuncio se obtienen $4$ manzanas $+ 2$
peras $= 10$ euros, pero el puesto cobra $9$: el sistema es incompatible,
rectas paralelas y conjunto de soluciones vacío. Veredicto: las dos
afirmaciones de precio no pueden ser ciertas a la vez; o se esconde un
descuento, o se esconde un error de cuentas.

\textbf{16.} En el origen: $0 > 3 \times 0 - 1 = -1$: cierto, así que
$(0,0)$ está en la región $y > 3x - 1$ (por encima de la recta). Cualquier
punto se comprueba sustituyéndolo: quedará por encima o por debajo según
la desigualdad sea cierta o falsa.

\textbf{17.} Vértices: $(0, 0)$; $(4, 0)$ (los ejes y $x + y = 4$);
$(0, 1)$ (el eje y $y = 2x + 1$); y el corte de $x + y = 4$ con
$y = 2x + 1$: $x = 1$, $y = 3$: $(1, 3)$.

\textbf{18.} $P(0,0) = 0$; $P(4, 0) = 12$; $P(0, 1) = 2$; $P(1, 3) = 9$.
\omterm{def:g10:functions:extrema}{Máximo} en el vértice $(4, 0)$: el plan $x = 4$, $y = 0$ rinde $12$; mandan
los vértices, como dejan ver las rectas paralelas $3x + 2y = c$ al barrer
la región.

\textbf{19.} $1.5x + 3 = 2x$ da $x = 6$, $y = 12$: a los seis kilómetros
los dos taxis cobran doce euros; el punto de equilibrio del problema de
fin de semana del volumen anterior sobre los dos termómetros, ahora
convertido en la solución única de un sistema.

\textbf{20.} Las soluciones de una \omterm{def:g10:algebra:equation}{ecuación} lineal dibujan una recta; un
sistema superpone dos rectas, y $ad - bc$ dice de un vistazo si se cortan
una vez (si es no nulo) o si caen en los casos paralelas-o-idénticas (si
es nulo). Si en lugar de eso apilamos inecuaciones, las soluciones forman
una región poligonal cuyos vértices llevan cualquier óptimo lineal. El
álgebra lineal generaliza el determinante a cualquier número de
incógnitas; la programación lineal industrializa los vértices.
\end{solution}
